\chapter{Functions and operators}\label{ch-procs}

A program of any size is built from parts that each solve one task and
that can be used again: functions. In Algol 68, a function is a value like
a number or an array. It has a type, it can be bound to a name, stored in
a variable, passed to another function and returned from one. The Report
calls a function a \rr{procedure}; this book uses both words, since
Algol 68 programs declare functions with the reserved word \code|proc|.

The second half of the chapter shows how a program defines operators of
its own. Operators are functions with a different calling syntax, and
they make it possible to extend the arithmetic of the language to new
types, such as fractions or vectors.

\section{Declaring functions}\label{sec-proc-decl}

A function declaration gives the function a name, its parameters with
their types, its result type and its body:

\example{ch05/declare}

The declaration \code|proc max = (int a, b) int: (a > b | a | b)| has the
same form as the constant declarations of Chapter~\ref{ch-values}: a name, an equals
sign and a value. The value is a \term{function literal}
(\rr{routine text}), the part \code|(int a, b) int: (a > b | a | b)|. It
lists the parameters in parentheses, then the result type, a colon and
the body, which is a single expression, usually a block. The type of
\code|max| is \code|proc (int, int) int|, the type of functions that take
two \code|int| values and return an \code|int|. Parameters of the same
type can share the type, as in \code|(int a, b)|; parameters of different
types are separated by commas: \code|(string s, int n)|.

The declaration \code|proc max = (int a, b) int: ...| is short for
\code|proc (int, int) int max = (int a, b) int: ...|, which names the
type in full. The short form requires a function literal on the
right-hand side. A function value computed in some other way, such as a
choice between two functions, needs the full type:
\code|proc (real) real f = (c | sin | cos)|. The short form with such a
value is rejected with a message about a missing symbol that does not
point to the cause.

A function whose result is of no interest has the result type
\code|void|, as \code|greet| shows. There is no \texttt{return}
statement: the result of a function is the value of its body, which is
the value of the last expression in the block. A function that must
finish early jumps to a label at the end of its body (Section~\ref{sec-jumps}),
which is rarely necessary since \code|if| and \code|case| yield values.

A call gives the arguments in parentheses after the name. The Report
calls the parameters \rr{formal parameters} and the arguments
\rr{actual parameters}.

\section{Parameters}\label{sec-params}

When a function is called, each parameter is bound to its argument
exactly as a constant is bound to its value in a declaration: the call
\code|max (3, 7)| behaves like the declarations \code|int a = 3, b = 7|
followed by the body. This explains the rules for parameters, which differ
from those of most languages.

A parameter is a constant. It cannot be assigned to:

\example{ch05/param-constant}

A function that needs a variable for its work declares one, for example
\code|int k := n|. A function that is to change a variable of its caller
takes a reference to that variable as its parameter:

\example{ch05/ref-param}

The parameters of \code|swap| are of type \code|ref int|. The call
\code|swap (x, y)| passes the references \code|x| and \code|y|, and the
assignments in the body go to the caller's variables. In C this would be
done with pointers, in C++ with reference parameters, in Pascal with
\texttt{var} parameters. Algol 68 needs no special kind of parameter for
it: the reference is an ordinary value, passed like any other.

Both kinds of parameter can be arrays. A parameter of type \code|[] real|
accepts an array of \code|real| with any bounds, and the function finds
them with \code|LWB| and \code|UPB|. A parameter of type \code|ref [] real|
accepts an array variable, whose elements the function can then change:

\example{ch05/array-param}

The second call of \code|sum| passes an array literal. Since a literal
has the bounds 1 to $n$, the function sees those bounds; for the variable
\code|w|, it sees the bounds 0 to 3.

A function can also return a reference to a cell that already exists,
such as an element of an array it receives by reference. Its call can
then stand on the left of an assignment, as in C++ a function that
returns a reference:

\example{ch05/max-elem}

\noindent The parameter must be a \code|ref [] real|. With
\code|[] real a|, the element \code|a[k]| would be a value, not a cell of
the caller's array, and could not be returned as a \code|ref real|.

\section{Functions without parameters}\label{sec-deproc}

A function may have no parameters. Its type is then \code|proc int| (or
\code|proc void|, and so on), and its literal has no parameter list:
\code|int: expression|. Such a function is called simply by writing its
name; no empty parentheses are needed, or even allowed:

\example{ch05/parameterless}

Writing \code|next| where an \code|int| is expected calls the function. This
is another implicit conversion, which the Report calls
\rr{deproceduring}. Where a function itself is expected, the name is not
a call: \code|newline| in a call of \code|print| is a function with a
parameter, passed as a value, and \code|print| calls it at the right
moment.

Each call of \code|next| changes the variable \code|seed|, so the two
calls give different results. In a list such as
\code|print ((next, next))| the order of the two calls would be
unspecified (Section~\ref{sec-collateral}); the program above therefore
calls the function in two separate declarations.

A parameterless function whose value is discarded is called as well,
which is how a function of type \code|proc void| is called as a
statement. (\code|newline| is not such a function: it has a file
parameter, and \code|newline;| alone does nothing; \code|print (newline)|
writes a line break.) There is one exception that is easy to overlook:

\example{ch05/discard-call}

\noindent prints \code|assigned| and then \code|hello| twice. The
assignment \code|p := hello| stores the function in the variable
\code|p| and does not call it, although its value, the variable
\code|p|, is discarded. A discarded assignment, a reference comparison,
an allocation, a type annotation, a literal and a format are discarded
as they are; any other expression that yields a function without parameters is
called.

A parameter of type \code|proc real| differs from one of type
\code|real| in when the argument is evaluated. A \code|real| parameter
receives a value, computed once before the call; a \code|proc real|
parameter receives a function, which is called at every use of the
parameter:

\example{ch05/by-name}

\noindent \code|twice_value| gets the value 1 and doubles it.
\code|twice_call| gets the function \code|count| and calls it twice,
which gives 2 and 3, and returns their sum, 5. Passing a literal such as \code|real: v[i] * w[i]| for a
\code|proc real| parameter gives the effect of Algol 60's call by name:
the expression is evaluated anew at each use, with the current value of
\code|i|. Algol 60 programs used this for summation functions, a trick
known as Jensen's device; in Algol 68 a parameter of type
\code|proc (int) real| does the same more plainly.

\section{Recursion}\label{sec-recursion}

A function may call itself. The classic examples are the factorial and
the Fibonacci numbers:

\example{ch05/recursion}

Two functions can also call each other. Since the scope of a declaration
is the whole block (Section~\ref{sec-scope}), \code|is_even| can mention
\code|is_odd| before \code|is_odd| is declared. This is allowed because
the call only happens later, when both declarations have been carried
out:

\example{ch05/mutual}

\section{Functions as values}\label{sec-proc-values}

A function is a value, so there are also variables that hold functions.
The declaration \code|proc (real) real f := sqrt| declares a variable for
functions from \code|real| to \code|real| and stores in it the square root
of the standard library. A later assignment can store another function,
for example a function literal written directly in the program:

\example{ch05/proc-values}

A function literal is what other languages call a \term{lambda} or an
anonymous function. Functions that take other functions as parameters are
called \term{higher-order functions}. A typical example computes an
integral numerically:

\example{ch05/higher-order}

The first call passes a declared function, the bell curve
$e^{-x^2}$, whose integral is the square root of $\pi$, about 1.7725; the second
passes a literal for $x^2$. The function in a call can itself be the
result of an expression: \code|(b | sin | cos) (x)| chooses the function
first and then calls it with \code|x|.

The body of a function can use the names of the blocks around it, not
only its parameters. A function can thus read and change variables of
the enclosing blocks:

\example{ch05/closure}

Such a function is called a \term{closure} in current terminology. Algol
68 closures have a restriction that the closures of Lisp, JavaScript or
Python do not have: a function may not be used after the block that
declares a name it uses has ended. A function that uses a variable, a
constant or a parameter of a block therefore cannot be returned from
that block or stored in a variable outside it.
Section~\ref{sec-proc-lifetimes} explains this rule, which is about the
lifetimes of values, and how the implementations check it.

\section{The standard functions}\label{sec-std-procs}

The standard library provides the usual mathematical functions for
\code|real| arguments: \code|sqrt|, \code|exp|, \code|ln| (the natural
logarithm), \code|sin|, \code|cos|, \code|tan|, \code|arcsin|,
\code|arccos| and \code|arctan|, together with the constant \code|pi|.
The parameterless function \code|random| returns a pseudo-random
\code|real| between 0 and 1. Chapter~\ref{ch-moretypes} describes the library in more
detail.

\section{Defining operators}\label{sec-op-decl}

An operator is declared like a function, with \code|op| instead of
\code|proc|. A binary operator has two parameters, a unary one has one.
The name of an operator is either a word in capitals, such as
\code|MAX|, or a symbol made of characters such as \code|+ - * / < > =|
and a few others. A symbol has one or two characters, and the second
must be one of \code|< > / = *|; either may be followed by \code|:=|
or \code|=:|. A unary operator may not begin with one of these five. Thus \code|//| and \code|+:=| can be declared, but
\code|++| cannot (``invalid operator tag''):

\example{ch05/operators}

A binary operator needs a precedence, declared with \code|prio|
(\rr{priority declaration}); without it, the compiler reports an error.
Here \code|MAX| gets the precedence 9, higher than that of \code|*|, so
\code|2 * 5 MAX 6| means \code|2 * (5 MAX 6)|. A unary operator needs no
precedence: like the standard unary operators, it binds more tightly
than every binary operator, so \code|SQUARE 4 + 1| is 17.

The same operator can be declared several times for different operand
types; the example declares \code|MAX| for two \code|int| and for two
\code|real| operands. This is \term{overloading}. The standard operators
are declared in the same way in the standard library, and a program can
add new meanings to them: \code|+| for vectors, \code|*| for matrices,
\code|<| for dates.

The declarations of one operator in the same block must differ enough
that the compiler can always tell them apart, even after the implicit
conversions allowed for operands. Two declarations clash if every
operand could be converted, by the conversions allowed for operands, to
the parameter types of both; for a binary operator this must hold for
both operands. Declarations for \code|(int, real)| and \code|(real, int)|
can therefore coexist. The Report calls clashing types \rr{firmly
related}. In practice the operand types must really differ, not merely
in whether an operand is a reference, and a union is related to each of
its members (Section~\ref{sec-expr}); the compiler reports a violation as
``firmly related''.

An operator can also be defined for arrays. The result of the following
\code|+| is a new array, which is copied out of the function like any
other value:

\example{ch05/vec-add}

The second operand is an array literal, which must be annotated with its
type: a literal is allowed only where its type is known, and the operand
of an operator is not such a place (Section~\ref{sec-casts}).

Since type names are only abbreviations, \code|mode Meters = real| does
not create a type of its own, and an operator for \code|Meters| would be
an operator for \code|real|. A record with a single field does create a
distinct type, the ``newtype'' of current languages. The record types
must have different field names, or they would be the same type:

\example{ch05/newtype}

A record literal needs at least two fields, so the functions
\code|meters| and \code|seconds| build the values.

\section{How an operator is chosen}\label{sec-op-identification}

When the compiler meets an operator, it chooses the declaration by the
types of the operands. Operands are in a firm context
(Section~\ref{sec-coercions}): a variable is dereferenced and a
parameterless function is called, but an \code|int| is not widened to a
\code|real|. An operator declared only for \code|real| operands therefore
does not accept an \code|int|:

\example{ch05/no-widening-op}

The standard operators avoid this problem by declaring every combination
of \code|int| and \code|real| operands.

The precedence belongs to the symbol, not to one of its declarations: all
the meanings of \code|+| in a block have the same precedence, and a
\code|prio| declaration changes it for all of them, the standard ones
included:

\example{ch05/prio-symbol}

\noindent prints 31 and 49. Inside the inner block, \code|3 + 4 * 7| is
\code|(3 + 4) * 7|. A program that wants the same
convenience must do the same, or call the operator with a \code|real|
operand. The Report calls the choice of the declaration the
\rr{identification} of the operator; current languages call it
\term{overload resolution}.

\section{A complete program}\label{sec-procs-complete}

The final program of this chapter defines a type of exact fractions and
the arithmetic operators for it. The function \code|rat| reduces a
fraction to lowest terms with the greatest common divisor and makes the
denominator positive. A new operator \code|//| builds a fraction from two
integers; it is given the precedence of \code|/|. The unary operator
\code|STR| converts a fraction into a string.

\example{ch05/rational}

Note how the standard symbols \code|+|, \code|-|, \code|*| and \code|/|
get a new meaning for the type \code|Rat| without losing their meaning
for numbers, and how the unary and binary \code|-| are declared
separately. Defining \code|+| does not define \code|+:=|: the program
declares it separately, with a reference to the variable as its left
operand and the variable as its result, like the standard \code|+:=|. The expression \code|STR - (6 // 8)| applies two unary
operators, the minus first.

\section*{Exercises}
\addcontentsline{toc}{section}{Exercises}

The \hyperref[ans-ch5]{answers} are at the end of the book.

\begin{enumerate}
\item What is the type of each of these functions?
  (a) \code|proc f = (real x, int n) real: x ** n|
  (b) \code|proc g = bool: random < 0.5|
  (c) \code|proc h = (ref [] int a) void: a[1] := 0|
\item Write a function \code|digit_sum| that returns the sum of the
  decimal digits of a non-negative \code|int|, once with a loop and once
  recursively.
\item Why does \code|proc inc = (int n) void: n +:= 1| not compile, and how
  must it be written to increase the caller's variable?
\item Write a function \code|apply_twice| that takes a function of type
  \code|proc (int) int| and an \code|int| and applies the function twice.
  Call it with a function literal that doubles its argument.
\item Declare an operator \code|ABOVE| so that \code|a ABOVE b| is
  \code|true| if \code|a| is greater than \code|b|, for \code|int|
  operands, with the precedence of \code|>|. What is the value of
  \code|1 + 2 ABOVE 2|?
\item With the operators of this chapter, what is \code|SQUARE 2 MAX 3|?
\item Why does the assignment \code|p := hello| not call the function
  \code|hello|, although its value is discarded?
\item In a block that contains \code|prio + = 7|, what is the value of
  \code|1 + 2 * 3|?
\item May \code|op MAX = (int a, real b) real| and
  \code|op MAX = (real a, int b) real| be declared in the same block?
\item Write an operator \code|*| that multiplies a \code|[] real| by a
  \code|real| and returns a new array with the same bounds.
\item A function \code|month = (ref int day) string| takes a day of the
  year, returns the name of its month and changes \code|day| to the day
  within that month. Why is \code|print ((month (d), " ", d))| wrong, and
  how is it corrected?
\item Write a function \code|sum = (int lo, hi, proc (int) real term) real|
  and use it to compute the sum of $1/k^2$ for $k$ from 1 to 100 and the scalar product of
  two arrays \code|v| and \code|w|. Why can the second call not simply
  pass \code|v[k] * w[k]|?
\end{enumerate}

\section*{Programming problems}
\addcontentsline{toc}{section}{Programming problems}

\begin{enumerate}
\item Write a function that returns the $n$-th prime number, and use it
  to print the first 20 primes.
\item Solve the Towers of Hanoi recursively and print the moves for four
  discs.
\item Write a function that finds a root of a function
  \code|proc (real) real| in an interval by bisection, and use it to
  compute the cube root of 2.
\item Define a type \code|Vector| with three \code|real| fields and the
  operators \code|+|, \code|-|, \code|*| (scalar product) and a unary
  \code|ABS| (length).
\item Extend the rational numbers of this chapter with the comparison
  operators \code|=| and \code|<|, and use them to sort an array of
  fractions.
\item \textbf{Merge sort.} Sort an array by recursive merge sort,
  passing the halves to the recursive calls as slices of a
  \code|ref [] int|.
\item \textbf{Permutations.} Print all permutations of the numbers 1 to
  $n$, and all ways of choosing $k$ of them, with a recursive procedure
  that fills an array position by position.
\item \textbf{Plotting.} Write a procedure that takes a function
  \code|proc (real) real|, an interval and a height, and plots the
  function with characters. Plot \code|sin| and a function literal.
\item \textbf{Horner's rule.} Write a function that evaluates a
  polynomial whose coefficients are given in an array with any lower
  bound, and count the multiplications compared with computing each power
  separately.
\item \textbf{The 91 function.} McCarthy's function is
  $f(n) = n - 10$ for $n > 100$ and $f(f(n + 11))$ otherwise. Compute it
  recursively for $n$ from 1 to 110, and find a simple formula for its
  values.
\item \textbf{Monte Carlo.} Estimate $\pi$ by choosing random points in
  a square with \code|random| and counting those inside the circle, for
  $10^3$, $10^4$ and $10^5$ points.
\item \textbf{Matrices.} Declare \code|+|, \code|*| and \code|*:=| for
  \code|[,] real|, checking that the bounds fit.
\item \textbf{Resistors.} Declare an operator \code|PAR| that combines
  two resistances in parallel, $ab/(a + b)$, choose a precedence for it
  relative to \code|+| (resistances in series), and add an assigning
  form \code|PARAB| in the style of \code|+:=|. Compute a network of
  several resistors with one formula.
\item \textbf{A maze.} Describe a maze by a function
  \code|blocked = (int x, y) bool| and find a way from the entrance to
  the exit with a recursive procedure that tries each direction, marks
  the cells it has visited, and backs out of dead ends. Print the path as the
  recursion returns, and explain why it comes out in reverse order.
\item \textbf{An evaluator.} Evaluate integer expressions with
  \code|+ - * /| and parentheses, given as a string, with mutually
  recursive functions for expressions, terms and factors that share a
  position variable.
\end{enumerate}
